\documentclass{jsarticle}
\usepackage{amsmath}
\usepackage{amssymb}
\begin{document}
\title{空間概念の量子化}
\author{Masaaki Yamaguchi}
\date{}
\maketitle

  あらまし

  始まりは、無限の時間の無が、場の理論として存在していた。
  この無が、オイラーの定数の大域的積分多様体が、
  基本群による、相対性から、オイラーの定数の多様体積分としての
  ガンマ関数に化けて、ゼウスとなり、このゼウスから、
  質量差エネルギーとして、母なるJones多項式ができた。
  このJones多項式から、各種の生成物として、
  宇宙と異次元が、D-braneを柱として、生み出された。
  このJones多項式から、中間子を媒介にして、素粒子が出来て、
  世界が切り開かれた。
  この世界は、Zeta関数を禅に対して、Beta関数が終わりに来るように、
  時間が２通りの流れとして、存在された。
  この時間の流れの中で、互いの相対性が消えて、
  空間概念の量子化としての、すべての共時性が最終目標にされた。

  登場人物は、

  大域的積分多様体のガンマ関数　父


  時間と空間　トキ子おばあちゃん


  Jones多項式　母と彩さん


  D-brane　山中崇史くん

  
  Zeta関数　無


  推移律の積分　先生たち


  アカシックレコードのヒント　姉たち


  著者　私




\end{document}
